% English references reproduced from the original edition.
% Citation-page links are generated dynamically by hyperref's pagebackref option.
\begingroup
\renewcommand{\bibname}{References}
\renewcommand{\bibpreamble}{%
  The page numbers in this book where a reference is cited are given after the
  reference with $\uparrow$. Each page number is a hyperlink back to the citation.\par\medskip
}
\phantomsection
\addcontentsline{toc}{chapter}{References}
\begin{thebibliography}{99}

\bibitem[Adler(2013)]{adler2013}
Adler, I. (2013). The equivalence of linear programs and zero-sum games. \emph{International Journal of Game Theory} 42(1), 165-177.

\bibitem[Albert et al.(2007)Albert, Nowakowski, and Wolfe]{albert2007}
Albert, M. H., R. J. Nowakowski, and D. Wolfe (2007). \emph{Lessons in Play: An Introduction to Combinatorial Game Theory}. AK Peters, Wellesley, MA.

\bibitem[Allais(1953)]{allais1953}
Allais, M. (1953). Le comportement de l'homme rationnel devant le risque: critique des postulats et axiomes de l'école américaine [with English summary]. \emph{Econometrica} 21(4), 503-546.

\bibitem[Aumann(1974)]{aumann1974}
Aumann, R. J. (1974). Subjectivity and correlation in randomized strategies. \emph{Journal of Mathematical Economics} 1(1), 67-96.

\bibitem[Aumann(1987)]{aumann1987}
Aumann, R. J. (1987). Correlated equilibrium as an expression of Bayesian rationality. \emph{Econometrica} 55(1), 1-18.

\bibitem[Aumann and Brandenburger(1995)]{aumann1995}
Aumann, R. and A. Brandenburger (1995). Epistemic conditions for Nash equilibrium. \emph{Econometrica} 63(5), 1161-1180.

\bibitem[Aumann and Dreze(2008)]{aumann2008}
Aumann, R. J. and J. H. Dreze (2008). Rational expectations in games. \emph{American Economic Review} 98(1), 72-86.

\bibitem[Avenhaus and Canty(1996)]{avenhaus1996}
Avenhaus, R. and M. J. Canty (1996). \emph{Compliance Quantified: An Introduction to Data Verification}. Cambridge University Press, Cambridge, UK.

\bibitem[Avis et al.(2010)Avis, Rosenberg, Savani, and von Stengel]{avis2010}
Avis, D., G. D. Rosenberg, R. Savani, and B. von Stengel (2010). Enumeration of Nash equilibria for two-player games. \emph{Economic Theory} 42(1), 9-37.

\bibitem[Berlekamp et al.(2001-2004)Berlekamp, Conway, and Guy]{berlekamp2001}
Berlekamp, E. R., J. H. Conway, and R. K. Guy (2001-2004). \emph{Winning Ways for Your Mathematical Plays}, 2nd ed., Volumes 1-4. AK Peters, Wellesley, MA. (First edition 1982).

\bibitem[Bernoulli(1738)]{bernoulli1738}
Bernoulli, D. (1738). Specimen theoriae novae de mensura sortis. \emph{Commentarii Academiae Scientiarum Imperiales Petropolitanae} 5, 175-192. Translated as: Exposition of a new theory on the measurement of risk, \emph{Econometrica} 22(1), 23-36, 1954.

\bibitem[Bewersdorff(2005)]{bewersdorff2005}
Bewersdorff, J. (2005). \emph{Luck, Logic, and White Lies: The Mathematics of Games}. AK Peters, Wellesley, MA.

\bibitem[Biggs et al.(1976)Biggs, Lloyd, and Wilson]{biggs1976}
Biggs, N., E. K. Lloyd, and R. J. Wilson (1976). \emph{Graph Theory 1736-1936}. Clarendon Press, Oxford.

\bibitem[Binmore(1992)]{binmore1992}
Binmore, K. (1992). \emph{Fun and Games}. D. C. Heath, Lexington, MA.

\bibitem[Binmore(2007)]{binmore2007}
Binmore, K. (2007). \emph{Playing for Real}. Oxford University Press, New York.

\bibitem[Bouton(1901)]{bouton1901}
Bouton, C. L. (1901). Nim, a game with a complete mathematical theory. \emph{Annals of Mathematics} 3(1/4), 35-39.

\bibitem[Braess(1968)]{braess1968}
Braess, D. (1968). Über ein Paradoxon aus der Verkehrsplanung. \emph{Unternehmensforschung} 12(1), 258-268. Translated in Braess, Nagurney, and Wakolbinger (2005).

\bibitem[Braess et al.(2005)Braess, Nagurney, and Wakolbinger]{braess2005}
Braess, D., A. Nagurney, and T. Wakolbinger (2005). On a paradox of traffic planning. \emph{Transportation Science} 39(4), 446-450.

\bibitem[Brouwer(1911)]{brouwer1911}
Brouwer, L. E. J. (1911). Über Abbildung von Mannigfaltigkeiten. \emph{Mathematische Annalen} 71(1), 97-115.

\bibitem[Brouwer(1976)]{brouwer1976}
Brouwer, L. E. J. (1976). \emph{Collected Works 2: Geometry, Analysis, Topology and Mechanics}. North-Holland, Amsterdam. Edited by H. Freudenthal.

\bibitem[Bryant(1990)]{bryant1990}
Bryant, V. (1990). \emph{Yet Another Introduction to Analysis}. Cambridge University Press, Cambridge, UK.

\bibitem[Chen et al.(2009)Chen, Deng, and Teng]{chen2009}
Chen, X., X. Deng, and S.-H. Teng (2009). Settling the complexity of computing two-player Nash equilibria. \emph{Journal of the ACM} 56(3), Article 14.

\bibitem[Cheng et al.(2004)Cheng, Reeves, Vorobeychik, and Wellman]{cheng2004}
Cheng, S.-F., D. M. Reeves, Y. Vorobeychik, and M. P. Wellman (2004). Notes on equilibria in symmetric games. In: \emph{Sixth Workshop on Game Theoretic and Decision Theoretic Agents at the 3rd Conference on Autonomous Agents and Multi-Agent Systems (AAMAS '04)}. New York.

\bibitem[Chvátal(1983)]{chvatal1983}
Chvátal, V. (1983). \emph{Linear Programming}. W. H. Freeman, New York.

\bibitem[Cohen(1967)]{cohen1967}
Cohen, D. I. A. (1967). On the Sperner lemma. \emph{Journal of Combinatorial Theory} 2(4), 585-587.

\bibitem[Conway(1977)]{conway1977}
Conway, J. H. (1977). All games bright and beautiful. \emph{American Mathematical Monthly} 84(6), 417-434.

\bibitem[Conway(2001)]{conway2001}
Conway, J. H. (2001). \emph{On Numbers and Games}, 2nd ed. AK Peters, Natick, MA. (First edition 1976).

\bibitem[Cottle et al.(1992)Cottle, Pang, and Stone]{cottle1992}
Cottle, R. W., J.-S. Pang, and R. E. Stone (1992). \emph{The Linear Complementarity Problem}. Academic Press, San Diego, CA.

\bibitem[Cournot(1838)]{cournot1838}
Cournot, A. A. (1838). \emph{Recherches sur les principes mathématiques de la théorie des richesses}. Hachette, Paris. Translated as: \emph{Researches into the Mathematical Principles of Wealth}, A. M. Kelly, New York, 1960.

\bibitem[Dantzig(1963)]{dantzig1963}
Dantzig, G. B. (1963). \emph{Linear Programming and Extensions}. Princeton University Press, Princeton, NJ.

\bibitem[Dixit and Nalebuff(1991)]{dixit1991}
Dixit, A. K. and B. J. Nalebuff (1991). \emph{Thinking Strategically: The Competitive Edge in Business, Politics, and Everyday Life}. W. W. Norton, New York.

\bibitem[Euler(1741)]{euler1741}
Euler, L. (1741). Solutio problematis ad geometriam situs pertinentis. \emph{Commentarii Academiae Scientiarum Imperiales Petropolitanae} 8, 128-140 + Plate VIII. Translated as: The solution of a problem relating to the geometry of position (Biggs, Lloyd, and Wilson, 1976, 3-8).

\bibitem[Fishburn(1970)]{fishburn1970}
Fishburn, P. C. (1970). \emph{Utility Theory for Decision Making}. Wiley, New York.

\bibitem[Fraenkel(1996)]{fraenkel1996}
Fraenkel, A. S. (1996). Scenic trails ascending from sea-level Nim to alpine Chess. In: \emph{Games of No Chance}, edited by R. J. Nowakowski, volume 29 of MSRI Publications, 13-42. Cambridge University Press, Cambridge, UK.

\bibitem[Freudenthal(1942)]{freudenthal1942}
Freudenthal, H. (1942). Simplizialzerlegungen von beschränkter Flachheit. \emph{Annals of Mathematics} 43(3), 580-582.

\bibitem[Friedman(1998)]{friedman1998}
Friedman, D. (1998). Monty Hall's three doors: Construction and deconstruction of a choice anomaly. \emph{The American Economic Review} 88(4), 933-946.

\bibitem[Gale(1960)]{gale1960}
Gale, D. (1960). \emph{The Theory of Linear Economic Models}. McGraw-Hill, New York.

\bibitem[Gale(1974)]{gale1974}
Gale, D. (1974). A curious Nim-type game. \emph{The American Mathematical Monthly} 81(8), 876-879.

\bibitem[Gale and Neyman(1982)]{gale1982}
Gale, D. and A. Neyman (1982). Nim-type games. \emph{International Journal of Game Theory} 11(1), 17-20.

\bibitem[Gibbons(1992)]{gibbons1992}
Gibbons, R. (1992). \emph{Game Theory for Applied Economists}. Princeton University Press, Princeton, NJ.

\bibitem[Gilboa and Zemel(1989)]{gilboa1989}
Gilboa, I. and E. Zemel (1989). Nash and correlated equilibria: Some complexity considerations. \emph{Games and Economic Behavior} 1(1), 80-93.

\bibitem[Grundy(1939)]{grundy1939}
Grundy, P. M. (1939). Mathematics and games. \emph{Eureka} 2, 6-8.

\bibitem[Hart and Mas-Colell(2000)]{hart2000}
Hart, S. and A. Mas-Colell (2000). A simple adaptive procedure leading to correlated equilibrium. \emph{Econometrica} 68(5), 1127-1150.

\bibitem[Hart and Schmeidler(1989)]{hart1989}
Hart, S. and D. Schmeidler (1989). Existence of correlated equilibria. \emph{Mathematics of Operations Research} 14(1), 18-25.

\bibitem[Herstein and Milnor(1953)]{herstein1953}
Herstein, I. N. and J. Milnor (1953). An axiomatic approach to measurable utility. \emph{Econometrica} 21(2), 291-297.

\bibitem[Hong et al.(2007)Hong, Song, and Wu]{hong2007}
Hong, S.-K., I.-J. Song, and J. Wu (2007). Fengshui theory in urban landscape planning. \emph{Urban Ecosystems} 10(3), 221-237.

\bibitem[Huizinga(1949)]{huizinga1949}
Huizinga, J. (1949). \emph{Homo Ludens: A Study of the Play-Element in Culture}. Routledge, London.

\bibitem[Knaster et al.(1929)Knaster, Kuratowski, and Mazurkiewicz]{knaster1929}
Knaster, B., C. Kuratowski, and S. Mazurkiewicz (1929). Ein Beweis des Fixpunktsatzes für n-dimensionale Simplexe. \emph{Fundamenta Mathematicae} 14(1), 132-137.

\bibitem[Kuhn(1950)]{kuhn1950}
Kuhn, H. W. (1950). Simplified two-person poker. In: \emph{Contributions to the Theory of Games, Vol. I}, edited by H. W. Kuhn and A. W. Tucker, volume 24 of Annals of Mathematics Studies, 97-103. Princeton University Press, Princeton, NJ.

\bibitem[Kuhn(1953)]{kuhn1953}
Kuhn, H. W. (1953). Extensive games and the problem of information. In: \emph{Contributions to the Theory of Games, Vol. II}, edited by H. W. Kuhn and A. W. Tucker, volume 28 of Annals of Mathematics Studies, 193-216. Princeton University Press, Princeton, NJ. Reprinted in Kuhn (1997).

\bibitem[Kuhn(1960)]{kuhn1960}
Kuhn, H. W. (1960). Some combinatorial lemmas in topology. \emph{IBM Journal of Research and Development} 4(5), 518-524.

\bibitem[Kuhn(1969)]{kuhn1969}
Kuhn, H. W. (1969). Approximate search for fixed points. In: \emph{Computing Methods in Optimization Problems 2}, edited by L. A. Zadeh, L. W. Neustadt, and A. V. Balakrishnan, 199-211. Academic Press, New York.

\bibitem[Kuhn(1997)]{kuhn1997}
Kuhn, H. W. (ed.) (1997). \emph{Classics in Game Theory}. Princeton University Press, Princeton, NJ.

\bibitem[Lemke(1965)]{lemke1965}
Lemke, C. E. (1965). Bimatrix equilibrium points and mathematical programming. \emph{Management Science} 11(7), 681-689.

\bibitem[Lemke and Howson(1964)]{lemke1964}
Lemke, C. E. and J. T. Howson, Jr (1964). Equilibrium points of bimatrix games. \emph{Journal of the Society for Industrial and Applied Mathematics} 12(2), 413-423.

\bibitem[Lewis(2016)]{lewis2016}
Lewis, M. (2016). \emph{The Undoing Project: A Friendship That Changed Our Minds}. W. W. Norton, New York.

\bibitem[Liu(1996)]{liu1996}
Liu, L. A. (1996). The invariance of best reply correspondences in two-player games. Working paper. City University of Hong Kong, Faculty of Business, Department of Economics and Finance.

\bibitem[Loomis(1946)]{loomis1946}
Loomis, L. H. (1946). On a theorem of von Neumann. \emph{Proceedings of the National Academy of Sciences of the United States of America} 32(8), 213-215.

\bibitem[Luce and Raiffa(1957)]{luce1957}
Luce, R. D. and H. Raiffa (1957). \emph{Games and Decisions: Introduction and Critical Survey}. Wiley, New York.

\bibitem[Mangasarian(1964)]{mangasarian1964}
Mangasarian, O. L. (1964). Equilibrium points of bimatrix games. \emph{Journal of the Society for Industrial and Applied Mathematics} 12(4), 778-780.

\bibitem[Maschler et al.(2013)Maschler, Solan, and Zamir]{maschler2013}
Maschler, M., E. Solan, and S. Zamir (2013). \emph{Game Theory}. Cambridge University Press, Cambridge, UK.

\bibitem[Matoušek and Gärtner(2007)]{matousek2007}
Matoušek, J. and B. Gärtner (2007). \emph{Understanding and Using Linear Programming}. Springer, Berlin.

\bibitem[Monderer and Shapley(1996)]{monderer1996}
Monderer, D. and L. S. Shapley (1996). Potential games. \emph{Games and Economic Behavior} 14(1), 124-143.

\bibitem[Morgenstern(1973)]{morgenstern1973}
Morgenstern, O. (1973). Ingolf Ståhl: Bargaining theory (book review). \emph{The Swedish Journal of Economics} 75(4), 410-413.

\bibitem[Moscati(2016)]{moscati2016}
Moscati, I. (2016). Retrospectives: How economists came to accept expected utility theory: The case of Samuelson and Savage. \emph{Journal of Economic Perspectives} 30(2), 219-236.

\bibitem[Moulin and Vial(1978)]{moulin1978}
Moulin, H. and J.-P. Vial (1978). Strategically zero-sum games: The class of games whose completely mixed equilibria cannot be improved upon. \emph{International Journal of Game Theory} 7(3/4), 201-221.

\bibitem[Myerson(1991)]{myerson1991}
Myerson, R. B. (1991). \emph{Game Theory: Analysis of Conflict}. Harvard University Press, Cambridge, MA.

\bibitem[Nash(1950)]{nash1950}
Nash, J. F., Jr. (1950). The bargaining problem. \emph{Econometrica} 18(2), 155-162. Reprinted in Kuhn (1997).

\bibitem[Nash(1951)]{nash1951}
Nash, J. (1951). Non-cooperative games. \emph{Annals of Mathematics} 54(2), 286-295. Reprinted in Kuhn (1997).

\bibitem[Nau et al.(2004)Nau, Gomez Canovas, and Hansen]{nau2004}
Nau, R., S. Gomez Canovas, and P. Hansen (2004). On the geometry of Nash equilibria and correlated equilibria. \emph{International Journal of Game Theory} 32(4), 443-453.

\bibitem[Nau and McCardle(1990)]{nau1990}
Nau, R. F. and K. F. McCardle (1990). Coherent behavior in noncooperative games. \emph{Journal of Economic Theory} 50(2), 424-444.

\bibitem[Neyman(1997)]{neyman1997}
Neyman, A. (1997). Correlated equilibrium and potential games. \emph{International Journal of Game Theory} 26(2), 223-227.

\bibitem[Osborne(2004)]{osborne2004}
Osborne, M. J. (2004). \emph{An Introduction to Game Theory}. Oxford University Press, New York.

\bibitem[Osborne and Rubinstein(1994)]{osborne1994}
Osborne, M. J. and A. Rubinstein (1994). \emph{A Course in Game Theory}. MIT Press, Cambridge, MA.

\bibitem[Owen(1967)]{owen1967}
Owen, G. (1967). Communications to the editor: An elementary proof of the minimax theorem. \emph{Management Science} 13(9), 765-765.

\bibitem[Papadimitriou(1994)]{papadimitriou1994}
Papadimitriou, C. H. (1994). On the complexity of the parity argument and other inefficient proofs of existence. \emph{Journal of Computer and System Sciences} 48(3), 498-532.

\bibitem[Papadimitriou and Roughgarden(2008)]{papadimitriou2008}
Papadimitriou, C. H. and T. Roughgarden (2008). Computing correlated equilibria in multi-player games. \emph{Journal of the ACM} 55(3), Article 14.

\bibitem[Park(1999)]{park1999}
Park, S. (1999). Ninety years of the Brouwer fixed point theorem. \emph{Vietnam Journal of Mathematics} 27(3), 187-222.

\bibitem[Perea(2012)]{perea2012}
Perea, A. (2012). \emph{Epistemic Game Theory: Reasoning and Choice}. Cambridge University Press.

\bibitem[Pigou(1920)]{pigou1920}
Pigou, A. (1920). \emph{The Economics of Welfare}. McMillan, London.

\bibitem[Pratt(1964)]{pratt1964}
Pratt, J. W. (1964). Risk aversion in the small and in the large. \emph{Econometrica} 32(1/2), 122-136.

\bibitem[Robinson(1951)]{robinson1951}
Robinson, J. (1951). An iterative method of solving a game. \emph{Annals of Mathematics} 54(2), 296-301. Reprinted in Kuhn (1997).

\bibitem[Rosenthal(1973)]{rosenthal1973}
Rosenthal, R. W. (1973). A class of games possessing pure-strategy Nash equilibria. \emph{International Journal of Game Theory} 2(1), 65-67.

\bibitem[Roughgarden(2016)]{roughgarden2016}
Roughgarden, T. (2016). \emph{Twenty Lectures on Algorithmic Game Theory}. Cambridge University Press, Cambridge, UK.

\bibitem[Roughgarden and Tardos(2002)]{roughgarden2002}
Roughgarden, T. and É. Tardos (2002). How bad is selfish routing? \emph{Journal of the ACM} 49(2), 236-259.

\bibitem[Rubinstein(1982)]{rubinstein1982}
Rubinstein, A. (1982). Perfect equilibrium in a bargaining model. \emph{Econometrica} 50(1), 97-109.

\bibitem[Savani and von Stengel(2006)]{savani2006}
Savani, R. and B. von Stengel (2006). Hard-to-solve bimatrix games. \emph{Econometrica} 74(2), 397-429.

\bibitem[Savani and von Stengel(2016)]{savani2016}
Savani, R. and B. von Stengel (2016). Unit vector games. \emph{International Journal of Economic Theory} 12, 7-27.

\bibitem[Schrijver(1986)]{schrijver1986}
Schrijver, A. (1986). \emph{Theory of Linear and Integer Programming}. John Wiley \& Sons, Chichester, UK.

\bibitem[Schwalbe and Walker(2001)]{schwalbe2001}
Schwalbe, U. and P. Walker (2001). Zermelo and the early history of game theory. \emph{Games and Economic Behavior} 34(1), 123-137.

\bibitem[Selten(1975)]{selten1975}
Selten, R. (1975). Reexamination of the perfectness concept for equilibrium points in extensive games. \emph{International Journal of Game Theory} 4(1), 25-55. Reprinted in Kuhn (1997).

\bibitem[Selvin(1975)]{selvin1975}
Selvin, S. (1975). Letters to the editor: A problem in probability. \emph{The American Statistician} 29(1), 67.

\bibitem[Shapley(1963)]{shapley1963}
Shapley, L. S. (1963). Some topics in two-person games. Memorandum RM-3672-PR. The Rand Corporation, Santa Monica, CA.

\bibitem[Shapley(1974)]{shapley1974}
Shapley, L. S. (1974). A note on the Lemke-Howson algorithm. \emph{Mathematical Programming Study} 1: Pivoting and Extensions, 175-189.

\bibitem[Siegel(2013)]{siegel2013}
Siegel, A. N. (2013). \emph{Combinatorial Game Theory}, volume 146 of Graduate Studies in Mathematics. American Mathematical Society, Providence, RI.

\bibitem[Sperner(1928)]{sperner1928}
Sperner, E. (1928). Neuer Beweis für die Invarianz der Dimensionszahl und des Gebietes. \emph{Abhandlungen aus dem Mathematischen Seminar der Universität Hamburg} 6(1), 265-272.

\bibitem[Sprague(1935)]{sprague1935}
Sprague, R. (1935). Über mathematische Kampfspiele. \emph{Tohoku Mathematical Journal (First Series)} 41, 438-444.

\bibitem[Ståhl(1972)]{stahl1972}
Ståhl, I. (1972). \emph{Bargaining Theory}. The Economic Research Institute at the Stockholm School of Economics, Stockholm. Reviewed in Morgenstern (1973).

\bibitem[Szpiro(2010)]{szpiro2010}
Szpiro, G. G. (2010). \emph{Numbers Rule: The Vexing Mathematics of Democracy, from Plato to the Present}. Princeton University Press, Princeton, NJ.

\bibitem[Szpiro(2020)]{szpiro2020}
Szpiro, G. G. (2020). \emph{Risk, Choice and Uncertainty: Three Centuries of Economic Decision-Making}. Columbia University Press, New York.

\bibitem[Turocy and von Stengel(2002)]{turocy2002}
Turocy, T. L. and B. von Stengel (2002). Game theory. In: \emph{Encyclopedia of Information Systems}, volume 2, 403-420. Elsevier Science (USA).

\bibitem[Tversky(1969)]{tversky1969}
Tversky, A. (1969). Intransitivity of preferences. \emph{Psychological Review} 76(1), 31-48.

\bibitem[van Damme(1987)]{vandamme1987}
van Damme, E. (1987). \emph{Stability and Perfection of Nash Equilibria}. Springer, Berlin.

\bibitem[von Neumann(1928)]{vonneumann1928}
von Neumann, J. (1928). Zur Theorie der Gesellschaftsspiele. \emph{Mathematische Annalen} 100(1), 295-320. Translated as: On the theory of games of strategy. In: \emph{Contributions to the Theory of Games, Vol. IV}, eds. A. W. Tucker and R. D. Luce, Annals of Mathematics Studies vol. 40, Princeton University Press, 1959, pp. 13-42.

\bibitem[von Neumann and Fréchet(1953)]{vonneumann1953}
von Neumann, J. and M. Fréchet (1953). Communication on the Borel notes. \emph{Econometrica} 21(1), 124-127.

\bibitem[von Neumann and Morgenstern(1947)]{vonneumann1947}
von Neumann, J. and O. Morgenstern (1947). \emph{Theory of Games and Economic Behavior}, 2nd ed. Princeton University Press, Princeton, NJ.

\bibitem[von Stackelberg(1934)]{vonstackelberg1934}
von Stackelberg, H. (1934). \emph{Marktform und Gleichgewicht}. Springer, Berlin. Translated as: \emph{Market Structure and Equilibrium}, Springer, Heidelberg, 2011.

\bibitem[von Stengel(1996)]{vonstengel1996}
von Stengel, B. (1996). Efficient computation of behavior strategies. \emph{Games and Economic Behavior} 14(2), 220-246.

\bibitem[von Stengel(1999)]{vonstengel1999}
von Stengel, B. (1999). New maximal numbers of equilibria in bimatrix games. \emph{Discrete and Computational Geometry} 21(4), 557-568.

\bibitem[von Stengel(2002)]{vonstengel2002}
von Stengel, B. (2002). Computing equilibria for two-person games. In: \emph{Handbook of Game Theory with Economic Applications}, edited by R. J. Aumann and S. Hart, volume 3, 1723-1759. North-Holland, Amsterdam.

\bibitem[von Stengel(2010)]{vonstengel2010a}
von Stengel, B. (2010). Follower payoffs in symmetric duopoly games. \emph{Games and Economic Behavior} 69(2), 512-516.

\bibitem[von Stengel(2021)]{vonstengel2021}
von Stengel, B. (2021). Finding Nash equilibria of two-player games. Preprint arXiv:2102.04580.

\bibitem[von Stengel and Zamir(2010)]{vonstengel2010b}
von Stengel, B. and S. Zamir (2010). Leadership games with convex strategy sets. \emph{Games and Economic Behavior} 69(2), 446-457.

\bibitem[Wardrop(1952)]{wardrop1952}
Wardrop, J. G. (1952). Some theoretical aspects of road traffic research. \emph{Proceedings of the Institution of Civil Engineers} 1(3), 325-362.

\bibitem[Wythoff(1907)]{wythoff1907}
Wythoff, W. A. (1907). A modification of the game of Nim. \emph{Nieuw Archief voor Wiskunde} 7(2), 199-202.

\bibitem[Young(2004)]{young2004}
Young, H. P. (2004). \emph{Strategic Learning and Its Limits}. Oxford University Press, Oxford.

\bibitem[Zermelo(1913)]{zermelo1913}
Zermelo, E. (1913). Über eine Anwendung der Mengenlehre auf die Theorie des Schachspiels. In: \emph{Proc. Fifth Congress Mathematicians, Cambridge 1912}, 501-504. Cambridge University Press, Cambridge, UK. Translated in Schwalbe and Walker (2001).

\bibitem[Ziegler(1995)]{ziegler1995}
Ziegler, G. M. (1995). \emph{Lectures on Polytopes}, volume 152 of Graduate Text in Mathematics. Springer, New York.

\end{thebibliography}
\endgroup
